\thispagestyle{otFront}\phantomsection\pdfbookmark[0]{Parcours de travail}{bm-parcours}
\otPageTitre{Parcours de travail}{Des bases à la recherche, à ton rythme}
\otIntertitre{Choisir le bon niveau}
\begin{tabularx}{\linewidth}{@{}lY@{}}\toprule
Exercices 1 à 5 & Appliquer une notion et vérifier les bases.\\[2mm]
Exercices 6 à 10 & Installer les méthodes, changer de représentation.\\[2mm]
Exercices 11 à 15 & Relier les notions et rédiger une démarche.\\[2mm]
Questions flash & Entretenir les automatismes, quelques minutes à chaque séance.\\[2mm]
Défis D01 à D16 & Chercher plus longtemps, avec des indications séparées.\\[2mm]
Bilans A et B & Mobiliser plusieurs chapitres en temps limité.\\\bottomrule
\end{tabularx}\par\bigskip
\otIntertitre{Une méthode en quatre temps}
\otConseil{1}{Lire les objectifs et le cours ; reformuler la propriété à utiliser.}
\otConseil{2}{Chercher sans corrigé, en dessinant et en nommant les inconnues.}
\otConseil{3}{Comparer les démarches ; repérer précisément la première difficulté.}
\otConseil{4}{Refaire l’exercice quelques jours plus tard, puis expliquer la solution à voix haute.}
\par\medskip\otIntertitre{Un travail réparti sur l’année}
L’ordre des chapitres est un choix éditorial, adaptable par le professeur. Les automatismes, le raisonnement et Python doivent être réinvestis dès les premiers thèmes, même si leurs pages de ressources sont regroupées à la fin.\par\medskip
Une séance courte peut réunir un calcul mental, un exercice direct et une question d’explication. Une séance longue peut être réservée à un problème ou à un défi ; note alors ce que tu as essayé, même si la recherche n’aboutit pas.\par\medskip
\begin{savoirfaire}[Suivre ses progrès]Le tableau de bord suit les 240 exercices de chapitre. Pour les défis et les bilans, noter séparément : réussi seul, avec une indication, ou à reprendre. Une correction comprise devient un objectif de révision, pas une simple page à cocher.\end{savoirfaire}
\clearpage